\originalchapter{习题课二: 连续电荷分布}
题源为\href{https://ocw.mit.edu/courses/8-02-physics-ii-electricity-and-magnetism-spring-2007/resources/solving02/}{Problem Solving 2}, PDF第1--6页. 下列保留P1的两题、P2、P3及考试练习题, 答题页重复部分合并. 所有解答为AI独立解答.

\begin{exercise}
P1 Q1. 半径$R$、长$L\gg R$的圆柱壳仅侧面均匀带总电荷$Q$, 两端无电荷. (a)求面密度并核单位; (b)在远离横向尺度$R$处将它看作线电荷, 其线密度是多少并核单位?
\end{exercise}
\begin{solution}
(a)侧面积$2\pi RL$, 所以$\sigma=Q/(2\pi RL)$, 单位C/m$^2$. (b)将整个截面的电荷压到同一轴段, $\lambda=Q/L$, 单位C/m. 只要求距离远大于$R$并不等于可以忽略有限长度$L$; 若距物体远大于$L$, 又应看作点电荷.
\end{solution}

\begin{exercise}
P1 Q2. 半径$R$、长$L\gg R$的实心圆柱内均匀带总电荷$Q$. (a)求体密度并核单位; (b)作上述线电荷近似时线密度是多少并核单位?
\end{exercise}
\begin{solution}
(a)$\rho=Q/(\pi R^2L)$, 单位C/m$^3$. (b)$\lambda=Q/L$, 单位C/m. 与圆柱壳相同的总量每长度并不表示两者在内部的场相同.
\end{solution}

\begin{exercise}
P2. 长$\ell$均匀正电荷棒总电荷$Q$, 点$P$在棒轴上、距一端$a>0$. (a)写场的积分; (b)求积分; (c)保持$Q,a$固定令$\ell\to0$, 核点电荷极限. 取$P$在棒左侧, 正$x$指向棒.
\end{exercise}
\begin{solution}
(a)以距近端$s$定位源, $0<s<\ell$, $\dd q=(Q/\ell)\dd s$, 场指向$-\uv x$, 因而$E_x=-\kc(Q/\ell)\int_0^\ell\dd s/(a+s)^2$.
(b)积分为$1/a-1/(a+\ell)$, 得$\ee=-\kc Q\uv x/[a(a+\ell)]$.
(c)$\ell\to0$给$-\kc Q\uv x/a^2$. 若改为保持线密度固定, 总电荷会趋于零, 不是原问所要的点电荷极限.
\end{solution}

\begin{exercise}
P3. 长$L$棒沿$x$轴从$-L/2$到$L/2$, 均匀带电$Q$. 求中垂线点$(0,y)$处电场的积分表达式, 并计算结果.
\end{exercise}
\begin{solution}
源位移为$-x\uv x+y\uv y$, 相反$x$元的横向贡献抵消. $y\ne0$时
\[
E_y=\kc\frac QL\int_{-L/2}^{L/2}\frac{y\dd x}{(x^2+y^2)^{3/2}}
=\frac{\kc Q}{y\sqrt{y^2+(L/2)^2}}.
\]
原函数可取$x/[y\sqrt{x^2+y^2}]$. 当$|y|\gg L$时是点电荷场; 当$L\gg|y|$且$\lambda=Q/L$固定时为$2\kc\lambda/y$, 是无限长线的场. 棒上的理想连续线电荷点不属于本式的有限场区.
\end{solution}

\begin{exercise}
考试练习题. 半径$R$的上半圆弧位于$xy$平面, 圆心为$P=(0,0)$, 弧上均匀带总电荷$Q$. 角$\theta$自正$x$轴逆时针从0到$\pi$. (A)写$P$处$E_y$的积分; (B)取无穷远为零求$P$处电势.
\end{exercise}
\begin{solution}
$\lambda=Q/(\pi R)$, $\dd q=(Q/\pi)\dd\theta$. 小元场指向圆心, 所以
$E_y=-\kc Q/(\pi R^2)\int_0^\pi\sin\theta\dd\theta=-2\kc Q/(\pi R^2)$.
横向分量为零. 所有电荷到$P$距离为$R$, 电势是$V(P)=\kc Q/R$. 这里电势相加没有正负投影, 场则必须带方向.
\end{solution}
